% Assigned sections for #7, #22, #23 and #29; shared-loop contribution under #18.
% Insert under the existing algorithm section, after the BFS/DFS subsections.
\subsection{Tìm kiếm chi phí đồng nhất (Uniform-Cost Search, UCS)}
\label{sec:bloxorz-ucs}

\subsubsection{Biểu diễn và giao diện sử dụng chung}
Trạng thái \texttt{State(cubes, bridges)} chứa các ô khối đang chiếm và trạng thái
đóng/mở của từng cầu. Khóa chuẩn hóa \texttt{State.key()} bao gồm cả hai trường;
hai vị trí khối giống nhau nhưng cấu hình cầu khác nhau vẫn là hai trạng thái.
Sau khi tách, các lập phương được sắp theo vị trí để tránh tìm kiếm lặp do đổi nhãn.

Hành động \texttt{(cube\_index, direction)} chỉ định lập phương trong trạng thái
nguồn và một trong bốn hướng. Chỉ số này có thể đổi sau khi trạng thái kế tiếp
được chuẩn hóa. Chọn điều khiển không tạo hành động riêng và không tốn lượt.
\texttt{Board.successors} cung cấp các cạnh hợp lệ, bao gồm cầu, công tắc, ô dễ
vỡ, dịch chuyển và hợp nhất; bộ giải không cài lại luật trò chơi.

\subsubsection{Hàm chi phí không đồng nhất}
Gọi $O(s)$ là tập ô bị chiếm trong trạng thái $s$ và $F$ là tập ô dễ vỡ. Chi phí
của cạnh hợp lệ $(s,a,s')$ được định nghĩa:
\begin{equation}
  c(s,a,s')=1+\left|\bigl(O(s')\setminus O(s)\bigr)\cap F\right|.
  \label{eq:bloxorz-action-cost}
\end{equation}
Mỗi lần lăn có chi phí nền $1$; mỗi ô dễ vỡ mới tiếp xúc cộng thêm $1$.
Khối chiếm nhiều nhất hai ô nên chi phí thuộc $\{1,2,3\}$ và
$c_{\min}=1>0$. Chính sách nằm tại \path{bloxorz.core.costs};
\texttt{Board.step\_cost} chỉ gọi chính sách này. Điều chỉnh trọng số không đòi
hỏi sửa luật chuyển trạng thái hoặc vòng lặp tìm kiếm.

Phụ phí thể hiện ưu tiên hạn chế tiếp xúc bề mặt dễ vỡ, theo ví dụ của đề bài.
Đây là mục tiêu tối ưu được thiết kế, không phải xác suất rơi: trò chơi vẫn xác
định. Hệ số $1$ tạo đánh đổi dễ giải thích, trong đó tránh một tiếp xúc tương
đương giảm một lần lăn thường; không lựa chọn hệ số bằng cách tối ưu kết quả
thực nghiệm. Lập phương đứng yên trên ô dễ vỡ không làm tăng chi phí mọi bước
của lập phương còn lại. Đích dịch chuyển được xét như tiếp xúc mới; cầu và công
tắc không có phụ phí riêng.

Hợp đồng hành động hiện hành không có cạnh đổi điều khiển độc lập. Vì vậy,
không thêm cạnh chi phí $0$ hoặc $1$ cho việc chọn lập phương; mọi cạnh di
chuyển thực sự vẫn có chi phí ít nhất $c_{\min}$. Bộ giải từ chối chi phí không
hữu hạn hoặc nhỏ hơn cận này. Tổng chi phí $g$ là tổng trọng số cạnh, còn độ
dài lời giải là \texttt{len(path)}; hai đại lượng không được đánh đồng.

\subsubsection{Hàng đợi ưu tiên và xử lý trạng thái lặp}
UCS lấy nút có chi phí tích lũy $g$ nhỏ nhất. Bảng \texttt{best\_cost} ánh xạ
khóa trạng thái chuẩn hóa tới chi phí tốt nhất đã phát hiện. Khi tìm được đường
rẻ hơn, cập nhật bảng và đưa một nút mới vào hàng đợi ưu tiên, kể cả khi trạng
thái từng được mở rộng. Do \texttt{heapq} không có thao tác giảm khóa trực tiếp,
bản ghi cũ bị bỏ qua khi lấy ra nếu $g$ không còn bằng chi phí tốt nhất.

Khóa ưu tiên gồm chi phí, giá trị phá hòa và bộ đếm thứ tự tăng dần, bảo đảm
kết quả lặp lại và không so sánh các đối tượng trạng thái. Liên kết cha và hành
động nguồn được dùng để khôi phục đường đi. Kiểm tra đích diễn ra khi lấy nút
hợp lệ ra khỏi hàng đợi. Vì mọi chi phí dương, lần lấy ra đó cho nghiệm có
tổng chi phí tối ưu; không nhất thiết ít bước nhất. Với đồ thị hữu hạn và
không áp giới hạn, thuật toán kết thúc bằng nghiệm hoặc kết luận không có nghiệm.

\subsubsection{Phần hạ tầng tìm kiếm chung}
Vòng lặp tại \path{bloxorz.solvers.base} quản lý việc phát hiện/cải thiện trạng
thái, đường đi, bộ đếm và giới hạn. Nó hỗ trợ thứ tự vào trước ra trước (FIFO),
vào sau ra trước (LIFO) và độ ưu tiên chi phí. FIFO/LIFO được giữ làm điểm mở
rộng cho người phụ trách BFS/DFS; phần bàn giao này không thêm các mô-đun đó.

\texttt{generated} đếm nút gốc và mỗi lần thực sự đưa nút vào hàng đợi;
\texttt{expanded} chỉ tăng ngay trước khi sinh trạng thái kế tiếp, sau kiểm tra
đích và giới hạn nút; \texttt{frontier\_peak} tính cả bản ghi cũ còn trong hàng
đợi. Bỏ qua bản ghi cũ không tăng số nút mở rộng. Kết quả giải được, đạt giới
hạn và không có nghiệm đều trả đầy đủ \texttt{Counts} theo cùng giao diện.

Theo mã giả của hợp đồng, giới hạn thời gian được kiểm tra trước khi lấy nút
ra, còn giới hạn nút được kiểm tra sau điều kiện đích và trước khi tăng bộ đếm.
Đồng hồ trong bộ giải chỉ thực thi giới hạn; thời gian và bộ nhớ thực nghiệm
thuộc trách nhiệm của bộ công cụ đo đạc chung. Việc nạp và kiểm tra màn chơi
diễn ra trước lời gọi \texttt{solve}; đối tượng \texttt{Board} sử dụng chỉ mục
đã biên dịch, không quét lại lưới trong hàm khởi tạo.

\subsection{Tìm kiếm A* (A-star Search)}
\label{sec:bloxorz-astar}

A* sử dụng cùng chi phí tại Công thức~\eqref{eq:bloxorz-action-cost}, ưu tiên
$f(s)=g(s)+h(s)$, với $h$ là hàm ước lượng (heuristic). Cơ chế cải thiện đường
đi và mở lại trạng thái giống UCS, bảo đảm một trạng thái từng gặp vẫn có thể
nhận đường rẻ hơn. Các bộ giải giữ nguyên giao diện
\texttt{solve(stage, caps)} và kiểu kết quả đã đóng băng.

\subsubsection{Thiết kế hàm ước lượng}
Trên màn không có công tắc tách, đặt:
\begin{align}
  d(s)&=\min_{(r,c)\in O(s)}\bigl(|r-r_G|+|c-c_G|\bigr),\\
  h(s)&=\frac{c_{\min}\,d(s)}{2}.
  \label{eq:bloxorz-heuristic}
\end{align}
Khoảng cách được lấy từ ô bị chiếm gần đích nhất, không từ một tọa độ tham
chiếu có thể làm quá cận. Dữ liệu tĩnh được kiểm tra một lần để tạo hàm ước
lượng, không quét lưới trên từng nút.

\subsubsection{Không đánh giá quá cao và tính nhất quán}
\textbf{Không đánh giá quá cao (admissibility).}
Một lần lăn giảm $d$ nhiều nhất hai ô. Bỏ qua hướng khối, vật cản, cầu và ô dễ
vỡ chỉ làm bài toán dễ hơn, nên cần ít nhất $d(s)/2$ bước để tới đích. Mỗi bước
tốn ít nhất $c_{\min}$, suy ra $h(s)\leq h^*(s)$, với $h^*$ là chi phí tối ưu
còn lại. Lập luận dùng cận chi phí của chính mô hình không đồng nhất đang chạy.

\textbf{Nhất quán (consistency).}
Với mỗi cạnh không dịch chuyển, $d(s)\leq d(s')+2$, do đó:
\begin{equation}
  h(s)\leq h(s')+c_{\min}\leq h(s')+c(s,a,s').
\end{equation}
Tại đích, $h=0$. Mọi phụ phí đều không âm nên bất đẳng thức vẫn đúng trên
các bước có tiếp xúc ô dễ vỡ.

\subsubsection{Giới hạn do dịch chuyển và hướng khối}
Công tắc tách có thể dịch chuyển xa tùy ý, phá vỡ giới hạn hai ô. Vì vậy dùng
$h=0$ trên \emph{mọi trạng thái của màn có công tắc tách}, kể cả trước khi tách.
Phương án này không đánh giá quá cao và nhất quán, nhưng làm A* trở thành UCS.
Chỉ tắt hàm ước lượng sau khi khối đã tách là không đủ vì chính cạnh kích hoạt
dịch chuyển cũng có thể làm giảm khoảng cách rất lớn.

Không cộng tùy ý một đơn vị phạt khi khối nằm. Nếu khối nằm ngay cạnh đích và
có thể đứng lên đích trong một bước chi phí $1$, khoảng cách ô gần nhất là $1$,
nên $h=1/2$. Cộng thêm $1$ cho hướng nằm tạo ước lượng $3/2>1$, vượt chi phí
thật. Vì vậy phần cài đặt không sử dụng hạng phạt hướng chưa được chứng minh.

\subsubsection{Kiểm chứng và phạm vi kết luận}
Kiểm thử phát lại đường đi, so sánh chi phí UCS/A*, kiểm tra giới hạn, trường
hợp không có nghiệm, bản ghi cũ và mở lại nút đã mở rộng dưới một hàm ước lượng
không đánh giá quá cao nhưng không nhất quán. Trên ba màn nhỏ có đáp án tính
tay, chi phí còn lại được tính độc lập ở mọi trạng thái có đường tới đích và
đối chiếu với $h$. Kiểm tra mở rộng trên mười màn chính thức duyệt $2\,946$
trạng thái đạt được, tạo đồ thị cạnh ngược để tính chi phí tối ưu còn lại và
kiểm tra tính nhất quán trên mọi cạnh.

Ví dụ kiểm tra tính đúng trên L2-02: phép duyệt ít bước nhất chỉ nằm trong
kiểm thử tìm đường $46$ bước, chi phí $64$; UCS/A* chọn đường $49$ bước, chi
phí $58$. Đây là bằng chứng cho sự khác nhau giữa mục tiêu số bước và mục tiêu
chi phí, không phải số liệu đo thời gian hoặc bộ nhớ. Các chứng minh ở trên
áp dụng cho mô hình đã nêu; phép duyệt hữu hạn là kiểm chứng bổ sung, không
thay thế chứng minh cho mọi màn có thể được tạo.

Với hàm ước lượng nhất quán đang dùng, mỗi trạng thái được mở rộng nhiều nhất
một lần hợp lệ. Gọi $V,E$ là số trạng thái/cạnh đạt được, thời gian dùng hàng
đợi ưu tiên là $O((V+E)\log(V+E))$, bộ nhớ $O(V+E)$ kể cả bản ghi thay thế.
Không suy ra A* luôn nhanh hơn UCS: độ chặt của cận và chi phí tính hàm ước
lượng đều ảnh hưởng kết quả. Trên màn tách, hai thuật toán có cùng chiến lược
vì $h=0$. Đánh giá hiệu năng toàn nhóm chờ bộ công cụ đo đạc và phiên bản tích
hợp cuối; phần này không đưa số đo chưa được thực hiện.
